\documentclass[11pt]{article}
\usepackage[T1]{fontenc}
\usepackage[utf8]{inputenc}
\usepackage{lmodern}
\usepackage[margin=1in]{geometry}
\usepackage{amsmath,amssymb,amsthm,mathtools}
\usepackage{booktabs,longtable,array}
\usepackage{microtype}
\usepackage{xurl}
\usepackage[colorlinks=true,linkcolor=blue,citecolor=blue,urlcolor=blue]{hyperref}
\newtheorem{theorem}{Theorem}[section]
\newtheorem{proposition}[theorem]{Proposition}
\newtheorem{lemma}[theorem]{Lemma}
\newtheorem{corollary}[theorem]{Corollary}
\theoremstyle{definition}
\newtheorem{definition}[theorem]{Definition}
\newtheorem{remark}[theorem]{Remark}
\newcommand{\R}{\mathbb R}
\newcommand{\C}{\mathbb C}
\newcommand{\dd}{\,d}
\newcommand{\ii}{\mathrm i}
\newcommand{\Rea}{\operatorname{Re}}
\newcommand{\Ima}{\operatorname{Im}}
\newcommand{\wind}{\operatorname{wind}}
\newcommand{\ind}{\operatorname{ind}}
\newcommand{\erf}{\operatorname{erf}}
\newcommand{\erfc}{\operatorname{erfc}}
\newcommand{\sgn}{\operatorname{sgn}}
\DeclareMathOperator{\Arg}{Arg}
\title{Root-Following Heat Contours, Phase Concentration,\\
and the Cap at Infinity\\
\large Analytical reductions and a signed surplus research program}
\author{Drafted for Edward Baker\\
\small With substantial LLM assistance from gpt-6-astra (Codex), effort ultra}
\date{20 October 2026}
\begin{document}
\maketitle
\begin{abstract}
Following simple zeros of a backward heat flow removes the side phase
integrals of its physical jet exactly. The remaining cap is a precise
collision counter: between two roots born at an ordinary event, its
phase decrease is \((N+1)\pi\), where \(N\) counts interior cap zeros.
A strictly proved surplus over \(\pi\) would therefore force a later
event, subject to finite continuation and boundary hypotheses.

We develop this signed surplus as an analytical estimating problem.
For the quadratic model we compute exact endpoint-layer masses and
finite tail bounds; sending the cap to infinity retains its half-turn.
For the complete theta source we derive a large-positive-time saddle
limit \(H(t,z/L_t)/H(t,0)\to\cos z\), where \(L_t\sim\frac14\log t\).
Fixed-index real roots approach zero, while their physical jet phase
concentrates near an interior extremum. Explicit source-moment error
bounds distinguish this theta layer from the quadratic endpoint layers.
Positive rescaling of the derivative channel changes concentration but
preserves winding.

These calculations quantify the phase accounting; they do not supply
the missing theta surplus. In particular, an adjacent fixed-index
cosine cap has the ordinary decrease \(\pi\), and a nonvanishing homotopy
to it cannot produce extra winding. We formulate the remaining
event-conditioned analytical tasks and separate them from earlier
bounded numerical exclusions and failed energy comparisons.
No later theta collision or improvement of the global threshold is claimed.
\end{abstract}
\tableofcontents

\section{The question and its precise scope}
\label{sec:scope}
The proposed contradiction starts from a hypothetical multiple real zero
at the all-real threshold of the theta heat family. It seeks a second
event at a strictly later time. The useful simplification is to follow
the two simple roots born at the first event, making the side phase
integrals vanish. The difficulty is then to prove that the closing cap
has accumulated more phase than the original pair requires.

The present revision organizes the argument around that missing
\emph{signed surplus}. Exact contour identities, concentration
estimates and the large-time source asymptotic are established below.
The arithmetic inequality that would make the surplus strictly positive
remains open. This distinction determines both the narrative and the
new analytical research program.

\subsection{The complete theta source}
For \(u\ge0\), set
\begin{align}
 \lambda_n&=\pi n^2,&
 \phi_n(u)&=(2\lambda_n^2e^{9u}-3\lambda_ne^{5u})
                      e^{-\lambda_ne^{4u}},&
 \Phi(u)&=\sum_{n\ge1}\phi_n(u), \label{eq:source}\\
 H(t,z)&=\int_0^\infty e^{tu^2}\Phi(u)\cos(zu)\dd u. \label{eq:H}
\end{align}
Each \(\phi_n(u)>0\), since
\(\lambda_ne^{4u}\ge\pi>3/2\). On compact complex \(t,z\) sets, the
source and each differentiated integrand have an integrable
superexponential majorant. Thus \(H\) is jointly analytic, is real and
even in real \(x\), and satisfies
\begin{equation}
                         H_t=-H_{xx}.                 \label{eq:heat}
\end{equation}
All its spatial derivatives are bounded on the whole real axis for a
fixed compact time interval. The sign in \eqref{eq:heat} matters:
\(u(s,y)=H(T-s,a+y)\) satisfies the ordinary forward heat equation
\(u_s=u_{yy}\). Throughout, ``later'' means increasing physical parameter
\(t\), hence the backward-heat direction.

Here is a normalization check independent of conventions for \(H\).
Let \(\Theta(v)=\sum_{n\in\mathbb Z}e^{-\pi n^2v}\) and
\(\psi(u)=\sum_{n\ge1}e^u e^{-\lambda_ne^{4u}}\).
The Gaussian Fourier transform followed by periodization gives
\(\Theta(v)=v^{-1/2}\Theta(1/v)\).
Consequently \(P(u)=e^u\Theta(e^{4u})=e^u+2\psi(u)\) is even,
\(\psi'(0)=-1/2\), and direct differentiation gives
\begin{equation}
             \Phi=(P''-P)/16=(\psi''-\psi)/8.          \label{eq:reflection}
\end{equation}
In particular \(\Phi\) extends evenly to \(\R\).
Splitting the Gaussian Mellin integral at \(1\), using theta inversion
on its lower part, and substituting \(v=e^{4u}\) gives, for
\(s=(1+\ii x)/2\),
\[
 \pi^{-s/2}\Gamma(s/2)\zeta(s)
  =\frac1{s(s-1)}+8\int_0^\infty\psi(u)\cos(xu)\dd u.
\]
The identity first follows in its convergence region and then by
analytic continuation. Twice integrating \(\psi''\) by parts proves
\begin{equation}
 8H(0,x)=\xi((1+\ii x)/2),\qquad
 \xi(s)=\tfrac12s(s-1)\pi^{-s/2}\Gamma(s/2)\zeta(s).   \label{eq:xi}
\end{equation}



\subsection{Threshold, event and continuation assumptions}
Classical de Bruijn--Newman theory gives a finite real \(\Lambda\) such
that all zeros of \(H(t,\cdot)\) are real exactly for \(t\ge\Lambda\)
\cite{debruijn,newman}. In the normalization above RH is equivalent to
\(\Lambda\le0\), while Rodgers and Tao proved \(\Lambda\ge0\)
\cite{rodgerstao}. These are background inputs, not results of this work.

The contour question is conditional on an \emph{attained ordinary event}
\[
 T=\Lambda>0,\qquad H(T,a)=H_x(T,a)=0,\qquad
 h_2:=H_{xx}(T,a)\ne0.
\]
An actual positive threshold is not known, and no theorem here says that
every hypothetical positive threshold is attained at finite height or
has an ordinary collision. The earlier numerical envelope
\(0<T\le1/5,\ a\in[1001,1002]\) is only one conditional class; it is
not a coverage statement for all possible events. No uniform lower
bound on \(|h_2|\) is assumed.

At a multiple zero the local heat polynomial has nonreal roots just
before the event. Thus a strictly later multiple zero would contradict
all-realness above \(T=\Lambda\); Appendix~\ref{app:index} supplies the
local proof, including higher multiplicities. The finite contour still
requires actual simple root graphs, finite continuation, a zero-free jet
on its boundary and a regular terminal cap. An unproved continuation,
escape from a chosen chart, or failure of a tracker is not a collision.

\section{Root-following walls reduce winding to the cap}
\label{sec:rootwalls}
Write
\[
 J=H+\ii H_x,\qquad E=H^2+H_x^2,\qquad
 Q=H_x^2-HH_{xx}.
\]
On a boundary where \(J\ne0\),
\begin{equation}
 d\arg J=
 \frac{(H_xH_{xx}-HH_{xxx})\,dt+
             (HH_{xx}-H_x^2)\,dx}{E}
 =\frac{Q_x\,dt-Q\,dx}{E}.                       \label{eq:winding}
\end{equation}
We orient space--time by \(dt\wedge dx\): the upper time cap runs
from left to right. Its \emph{phase decrease} is the signed quantity
\[
 M_{\rm cap}:=-\Delta_{\rm cap}\arg J=\int_{\rm cap}Q/E\,dx.
\]
It is a positive measure only when \(Q\ge0\); estimates below state
explicitly when they concern net phase rather than total variation.

If \(x=w(t)\) is a simple root graph, implicit differentiation gives
\[
 w'=\frac{H_{xx}}{H_x}.
\]
Pulling \eqref{eq:winding} back to the wall gives
\[
 \frac{d}{dt}\arg J(t,w(t))
 =\frac{H_{xx}}{H_x}-w'=0.
\]
Equivalently \(J=\ii H_x\) has constant phase there. It is the full
tangential expression that cancels; its time and moving-space terms
need not vanish separately.

\begin{lemma}[Exact phase between simple roots]\label{lem:cap}
Let \(f\) be real \(C^2\) on \([u,v]\), with simple endpoint roots and
exactly \(N\) interior roots, all simple. Then
\[
 \Delta_{u\to v}\arg(f+\ii f')=-(N+1)\pi.
\]
\end{lemma}
\begin{proof}
Between consecutive roots \(f\) has one strict sign. If \(f>0\), its
jet stays in the right half-plane, starts on the positive imaginary
axis and ends on the negative imaginary axis, with increment \(-\pi\).
For a negative lobe multiply the jet by \(-1\). Sum over the \(N+1\)
lobes. No pointwise phase monotonicity is needed.
\end{proof}

\begin{proposition}[Finite root strip]\label{prop:strip}
Suppose the two roots born at \((T,a)\) continue as compact simple graphs
\(w_-(t)<w_+(t)\) for \(T+\eta\le t\le T+R\). Choose \(\eta>0\)
in the local splitting neighborhood, so the bottom roots are adjacent.
Assume the upper cap has only simple zeros, with \(N\) in its interior.
Then
\begin{equation}
 M_{\rm cap}=(N+1)\pi,\qquad N\ \hbox{even},\qquad
 W_{\rm strip}=-N/2.                              \label{eq:strip}
\end{equation}
If an enclosing lower piece adds only the original ordinary event,
the full physical winding is \(-1-N/2\).
\end{proposition}
\begin{proof}
The upper cap contributes \(-(N+1)\pi\), the reversed lower cap
contributes \(+\pi\), and both walls contribute zero. Division by
\(2\pi\) proves the winding formula and its integrality gives the parity.
At the original ordinary event, the Jacobian of \((H,H_x)\) is
\[
 \begin{pmatrix}-h_2&0\\-h_3&h_2\end{pmatrix},
 \qquad h_3=H_{xxx}(T,a),
\]
with determinant \(-h_2^2\), hence index \(-1\).
Index additivity gives the final assertion. The general event count
is proved in Appendix~\ref{app:index}.
\end{proof}

Consequently a complete signed estimate \(M_{\rm cap}>\pi\) would force
\(N\ge2\) and at least one later event. Assuming continued adjacency
would instead impose \(N=0\) from the start. For the contradiction route,
simple continuation is a legitimate wall premise; a global cap count
cannot also be assumed as an estimating input.

\subsection{The initial event cannot be dropped with the apex}
A contour through the event has \(J=0\) on its boundary, so its physical
winding is undefined. One may truncate above the event and use
Proposition~\ref{prop:strip}, enclose it by a lower connector, or use
the ordinary-event quotient
\[
 \mathcal L(\tau,y)=-h_2\tau+\ii(h_2y-h_3\tau),\qquad Z=J/\mathcal L.
\]
The real linear map \(\mathcal L\) is invertible, has index \(-1\),
and \(J-\mathcal L=O(\tau^2+y^2)\). Thus \(Z\to1\) locally and has
a nonzero continuous apex extension. Its neighborhood is event
dependent; no uniform curvature floor is hidden in this argument.
Longer remainder and apex estimates are in
\href{../notes/01_FINITE_CONTOUR_AND_SOURCE_CONTRACT_20261020.md}{Note 01}.

For the exact quadratic \(H=y^2-2\tau\), the root cap has phase
\(-\pi\). The lower connector \(\tau=y^2-\eta\), traversed from
\(\sqrt{2\eta}\) to \(-\sqrt{2\eta}\), has
\(J=2\eta-y^2+2\ii y\) and contributes another \(-\pi\) through the
right half-plane. Its contribution does not tend to zero as
\(\eta\downarrow0\). The full winding is exactly the original \(-1\).

\section{The strict signed surplus and its error budget}
\label{sec:surplus}
Partition a regular root cap into two disjoint endpoint intervals and
their complementary middle. Let \(M_{\rm ends}\) and \(M_{\rm mid}\)
be their signed phase decreases. Proposition~\ref{prop:strip} gives
\begin{equation}
 \Delta:=M_{\rm ends}+M_{\rm mid}-\pi=N\pi
                 \in 2\pi\mathbb Z_{\ge0}.       \label{eq:surplus}
\end{equation}
\begin{corollary}[A sufficient strict budget]\label{cor:budget}
Under the finite root-strip hypotheses, complete independent bounds
\[
 M_{\rm ends}\ge\pi-\eta_R,\qquad M_{\rm mid}>\eta_R
\]
force a later event. More generally any complete lower bound
\(M_{\rm ends}+M_{\rm mid}>\pi\) suffices.
\end{corollary}
Indeed the next allowed value after \(\pi\) is \(3\pi\).
The estimates must include signed endpoint deficits, all source
remainders and the actual wall/cut locations. Concentration alone
does not provide either strict inequality.

Equation~\eqref{eq:surplus} is an exact reformulation of the extra-root
target under the stated hypotheses, not a new estimate or a weaker
intermediate theorem. Its practical use is to specify where a source
calculation must gain a strict margin. One may use exact endpoint
arguments rather than integrate the endpoint pieces.

\subsection{A scalar description and the local sign}
The born walls inherit slope signs
\(-\sgn h_2\) on the left and \(+\sgn h_2\) on the right.
Hence \(h_2H<0\) immediately inside both cap endpoints. A point \(b\)
strictly between the walls with
\begin{equation}
                  h_2H(T+R,b)>0                  \label{eq:sign-target}
\end{equation}
would force an interior root on each side. Conversely extra simple
interior roots produce such a sign interval. This is equivalent cap
information; it does not supply an arithmetic sign by itself.
Writing \(h_j=\partial_x^jH(T,a)\), the heat equation gives
\[
 h_2H(T+R,a)=-h_2^2R+\tfrac12h_2h_4R^2+O(R^3)<0
\]
for sufficiently small \(R>0\). The center is inside the born pair
at these small times. Better local quadratic approximation therefore
confirms the original lobe, rather than generating another one.

\subsection{What a phase homotopy can and cannot prove}
For complex cap approximants \(P,M\) with
\(|P-M|\le e<|M|\), the quotient \(P/M\) lies in a disk centered
at 1 that avoids zero. A continuous argument difference satisfies
\[
 |\arg(P/M)|\le\arcsin(e/|M|).
\]
Its two endpoint errors bound the difference of the net phase increments.
This retains phase correlation across the interval.
On a complete nonzero closed contour, a nonvanishing homotopy to an
ordinary one-event model certifies that same degree. It cannot
simultaneously force an extra event. A useful strict-surplus estimate
must use source information that goes beyond such a one-lobe enclosure.

\section{Quantitative endpoint concentration in the quadratic}
\label{sec:quadratic}
Take \(f_R(y)=y^2-r^2\), \(r=\sqrt{2R}\), on \([-r,r]\).
Here
\[
 Q=2(y^2+r^2),\qquad E=(y^2-r^2)^2+4y^2,\qquad
 d\mu_R=-d\arg(f_R+\ii f_R')=\frac{Q}{E}\,dy.
\]
This is a positive measure of total mass \(\pi\).
With \(z=y/r,\ \epsilon=2/r\),
\begin{equation}
 J/r^2=z^2-1+\ii\epsilon z,\qquad
 \frac{d\arg J}{dz}
 =-\frac{\epsilon(1+z^2)}{(1-z^2)^2+\epsilon^2z^2}. \label{eq:quad-density}
\end{equation}
A continuous lift is
\[
 \varphi_R(z)=-\pi-\arctan\frac{\epsilon z}{1-z^2},
 \quad -1<z<1,
\]
with endpoint values \(-\pi/2,-3\pi/2\).

\begin{proposition}[Finite endpoint-layer bounds]\label{prop:layers}
For \(0<\delta<1\), the middle interval
\([-1+\delta,1-\delta]\) has phase mass
\begin{equation}
 M_{\rm mid}
 =2\arctan\frac{\epsilon(1-\delta)}{\delta(2-\delta)}
 \le\frac{2}{r\delta}.                            \label{eq:middle-bound}
\end{equation}
Each endpoint interval has mass \(\pi/2-M_{\rm mid}/2\).
For physical endpoint windows of width \(0<V<r\), each has mass
\begin{equation}
 M_{\rm end}
 =\arctan\left(V+\frac{V^2}{2(r-V)}\right),\qquad
 0\le M_{\rm outside}\le\frac2V.                  \label{eq:physical-window}
\end{equation}
Moreover
\[
 0\le M_{\rm end}-\arctan V
 \le\frac{V^2}{2(r-V)(1+V^2)}.
\]
\end{proposition}
\begin{proof}
Evaluate the lift at the middle endpoints to get
\eqref{eq:middle-bound}; use \(\arctan x\le x\) and
\(2(1-\delta)\le2-\delta\).
Set \(\delta=V/r\). The complementary half-mass equals
\[
 \frac\pi2-\arctan\frac{2(r-V)}{V(2r-V)}
 =\arctan\left(V+\frac{V^2}{2(r-V)}\right).
\]
The outside mass is twice the arctangent of the reciprocal of this
last argument, at most \(2/V\). Finally
\(\arctan(V+B)-\arctan V
=\arctan[B/(1+V(V+B))]\le B/(1+V^2)\), with
\(B=V^2/[2(r-V)]\).
\end{proof}

Thus \(V\ge2/\eta\) and \(R>V^2/2\) put all but \(\eta\) radians
in the two physical endpoint windows. Their widths remain \(O(1)\);
their relative widths vanish because the cap grows like \(\sqrt R\).
Equivalently, in the \(z\) coordinate,
\[
 d\mu_R\rightharpoonup
       \tfrac\pi2\delta_{-1}+\tfrac\pi2\delta_1.
\]
This follows from symmetry, total mass \(\pi\), and
\eqref{eq:middle-bound}. Pointwise interior vanishing is not convergence
of the total phase to zero. The limiting normalized jet \(z^2-1\)
vanishes at the endpoints.

The strict budget is exactly saturated:
\[
                  M_{\rm mid}=\pi-M_{\rm ends}.
\]
Keeping the small middle phase while replacing the endpoints by their
limiting mass \(\pi\) would introduce a false surplus.

\subsection{A reusable root-window estimate}
\begin{lemma}[Relative curvature controls endpoint phase]\label{lem:curvature}
At a simple left root \(u\), put \(s=f'(u)\ne0\). Suppose
\[
 |f''(u+d)|\le K|s|\quad(0\le d\le V),\qquad \rho=KV<1.
\]
Then the window contains no further root, and its net phase decrease
obeys
\[
 \arctan\frac{V(1-\rho/2)}{1+\rho}
 \le M_{\rm end}\le
 \arctan\frac{V(1+\rho/2)}{1-\rho}.
\]
The inward right-endpoint formula is identical.
\end{lemma}
\begin{proof}
Taylor's integral remainder gives
\(f(u+d)/s=d+e_0,\ f'(u+d)/s=1+e_1\),
with \(|e_0|\le Kd^2/2,\ |e_1|\le Kd\).
These normalized components are positive inside the window.
The phase decrease is exactly
\(\arctan[f(u+V)/f'(u+V)]\); bound that positive ratio.
\end{proof}
Capturing nearly \(\pi/2\) needs both large \(V\) in the jet's unit
derivative scale and small \(KV\). Simplicity alone supplies neither.
The quadratic has \(K=1/r\). A theta application needs its own relative
curvature estimate and enough separation for the two windows.

\section{What changes when the cap is sent to infinity?}
\label{sec:infinity}
Start with the complete finite contours of Proposition~\ref{prop:strip}.
The endpoint unit phases remain opposite imaginary directions.
A continuous cap joining them cannot have constant phase, even at a
limiting boundary. If they stay adjacent for every finite \(R\),
\(M_{\rm cap}(R)=\pi\) identically.

When the finite interior events are isolated, the root-strip count
increases by twice their collision charges. A finite limiting count
is eventually constant. A limiting strict excess would therefore
already occur at a finite regular cap. An unbounded count gives
unbounded phase, not a finite winding at infinity. A genuine limiting
argument must control the full phase and identify the same walls;
an amplitude limit or pointwise phase-density limit is insufficient.

\subsection{A continuous limiting cap can retain its half-turn}
The positive cosine-source family
\[
 H(t,x)=1+e^{t-T_0}\cos x
\]
has actual all-real threshold \(T_0\). Indeed \(\cos x=-e^{T_0-t}\)
has only real solutions for \(t\ge T_0\), while for \(t<T_0\) it has
nonreal conjugate solutions at each odd multiple of \(\pi\).
Centering at \(y=x-\pi,\ \tau=t-T_0\), its born walls are
\(\pm\alpha_\tau\), \(\alpha_\tau=\arccos(e^{-\tau})\).
Between them \(H=1-e^\tau\cos y<0\), with no later collision.
On the cap \(y=\alpha_R z\),
\[
 J/e^R=e^{-R}-\cos(\alpha_R z)+\ii\sin(\alpha_R z)
       \longrightarrow-e^{-\ii\pi z/2},\qquad -1\le z\le1.
\]
The convergence is uniform to a nonzero curve, whose phase change is
still \(-\pi\). The quadratic loses a nonvanishing endpoint limit;
this atomic control has a good limit and still retains the cap.
Neither supports dropping it.

\subsection{The detector does see a genuine later event}
For
\[
 F(s,y)=y^4-12sy^2+12s^2-3
\]
the two ordinary events are \((-1/2,0)\) and \((1/2,0)\).
The outer roots born at the first event are
\[
 w_\pm(s)=\pm\sqrt{6s+\sqrt{24s^2+3}}.
\]
They remain simple. For \(s>1/2\) there are two interior simple roots,
giving cap decrease \(3\pi\) and full physical winding \(-2\).
The first event is not the all-real threshold: all four roots are real
exactly for \(s\ge1/2\). The example verifies the signed detector
without contradicting future simplicity above a true threshold.

% Prepared for Edward Baker, 20 October 2026, with substantial LLM assistance.
% Model gpt-6-astra; reasoning effort ultra, verified by parent metadata.
% Separate same-model analytical draft; no specialist review.
\section{The complete theta profile and its interior phase layer}
\label{asym:section}
This section concerns $t\to+\infty$, with fixed spatial indices and
fixed rescaled compact sets.  It gives a complete-source asymptotic,
not a continuation theorem for previously selected walls or a collision
criterion.  The heat convention remains $H_t=-H_{xx}$.

\subsection{A positive spectral law and its large stationary point}
Use the complete normalization
\begin{equation}\label{asym:source}
 H(t,x)=\int_0^\infty e^{tu^2}\Phi(u)\cos(xu)\dd u,\qquad
 \Phi(u)=\sum_{n\ge1}(2\pi^2n^4e^{9u}-3\pi n^2e^{5u})
                         e^{-\pi n^2e^{4u}}.
\end{equation}
Every summand is positive for $u\ge0$.  Write
\begin{equation}\label{asym:factor}
 e^{tu^2}\Phi(u)=2\pi^2e^{q_t(u)}b(u),\qquad
 q_t(u)=tu^2+9u-\pi e^{4u}.
\end{equation}
Retaining the first summand for a lower bound and dropping the negative
terms for an upper bound gives
\begin{equation}\label{asym:b-bounds}
 0<c_b:=1-\frac3{2\pi}\le b(u)\le
 C_b:=\sum_{n\ge1}n^4e^{-\pi(n^2-1)}<\infty,
 \qquad b(u)\longrightarrow1.
\end{equation}
Indeed the exact series for $b$ has terms
$[n^4-3n^2e^{-4u}/(2\pi)]e^{-\pi(n^2-1)e^{4u}}$.
The displayed summable majorant also justifies the limit.

For sufficiently large $t$, let $L=L_t$ be the large solution of
\begin{equation}\label{asym:saddle}
 4\pi e^{4L}=2tL+9,
 \qquad \kappa=\kappa_t:=16\pi e^{4L}-2t=8tL-2t+36.
\end{equation}
This choice is unambiguous: $q_t''=2t-16\pi e^{4u}$ strictly decreases,
$q_t'(0)=9-4\pi<0$, and its maximum is positive for large $t$.
Thus $q_t'$ has a small zero giving a local minimum and one large zero
giving the selected maximum.  The large zero tends to infinity.
Taking logarithms in \eqref{asym:saddle} first gives
$\log t=4L-\log L+\log(2\pi)+o(1)$ and hence
\begin{equation}\label{asym:L-expansion}
 L=\frac{\log t+\log\log t-\log(8\pi)+o(1)}4,
 \qquad \kappa\sim8tL.
\end{equation}

\begin{proposition}[Normalization and spectral concentration]
\label{asym:concentration}
Let $U$ have the probability law
\[
 d\mu_t(u)=\frac{e^{tu^2}\Phi(u)\dd u}{H(t,0)},\qquad u\ge0.
\]
Then
\begin{equation}\label{asym:normalization}
 H(t,0)\sim2\pi^2e^{q_t(L)}\sqrt{\frac{2\pi}{\kappa}}.
\end{equation}
The variable $\sqrt\kappa(U-L)$ converges to a standard normal variable,
with convergence of every fixed absolute moment.  In particular, setting
$m_1=\mathbb E|U-L|$ and $m_2=\mathbb E(U-L)^2$, one has
\begin{equation}\label{asym:moments}
 m_1\sim\sqrt{\frac2\pi}\,\kappa^{-1/2},\qquad
 m_2\sim\kappa^{-1}.
\end{equation}
\end{proposition}
\begin{proof}
Set $A_t=\pi e^{4L}$.  Cancellation of the linear term gives exactly
\begin{equation}\label{asym:exponent-difference}
 q_t(L+v)-q_t(L)=tv^2-A_t(e^{4v}-1-4v).
\end{equation}
For $v\ge0$, the right side is at most $-\kappa v^2/2$.
Fix $\delta>0$.  On $-\delta\le v\le0$, Taylor's integral remainder
bounds $e^{4v}-1-4v$ below by $8e^{-4\delta}v^2$.
Since $t/\kappa\to0$, for all sufficiently large $t$ this gives
\begin{equation}\label{asym:local-tail}
 q_t(L+v)-q_t(L)\le-c_\delta\kappa v^2,
 \qquad c_\delta=e^{-4\delta}/4.
\end{equation}
On $[0,L-\delta]$, the maximum of $q_t$ is at an endpoint, by its
one-minimum/one-maximum shape.  Eventually
$q_t(L-\delta)>q_t(0)$, since the former is asymptotic to $tL^2$.
Consequently its unnormalized mass is at most
\begin{equation}\label{asym:far-left}
 2\pi^2 C_b L\exp\{q_t(L)-c_\delta\kappa\delta^2\}.
\end{equation}
For $v=w/\sqrt\kappa$, \eqref{asym:exponent-difference} equals
$-w^2/2+O(|w|^3/\sqrt\kappa)$ on each fixed bounded $w$-interval;
here $A_t/\kappa\to1/16$.  Also $b(L+w/\sqrt\kappa)\to1$.
The Gaussian bounds on the remaining two regions and
\eqref{asym:far-left} prove \eqref{asym:normalization} by dominated
convergence.  After normalization the far-left probability is
$O(L\sqrt\kappa e^{-c_\delta\kappa\delta^2})$.
For an absolute moment of order $p$, its contribution is at most a
constant times $(L\sqrt\kappa)^{p+1}e^{-c_\delta\kappa\delta^2}$,
which tends to zero.  Polynomially weighted Gaussian domination handles
the other regions and proves all the stated moment limits.
\end{proof}

\subsection{Cosine profile, fixed-index roots, and the physical jet}
Define the entire rescaled profile
\begin{equation}\label{asym:F}
 F_t(z)=\frac{H(t,z/L)}{H(t,0)}=\mathbb E\cos(zU/L).
\end{equation}
Proposition~\ref{asym:concentration} implies
\begin{equation}\label{asym:cosine-limit}
 F_t(z)\longrightarrow\cos z
\end{equation}
locally uniformly on $\C$, with every fixed derivative.  To check
complex arguments, on $|z|\le Z$ use the factor $e^{ZU/L}$.
On the far-left region it is at most $e^Z$; on the right, after the
change to $w$, its extra factor is $e^{Zw/(L\sqrt\kappa)}$,
absorbed by the Gaussian bound.  Fixed derivative powers of $U/L$
are handled the same way.  Thus no oscillatory tail is discarded.

For completeness the real estimates needed below are explicit:
\begin{align}
 |F_t(z)-\cos z|&\le\varepsilon_0:=Zm_1/L,\nonumber\\
 |F_t'(z)+\sin z|&\le\varepsilon_1:=(1+Z)m_1/L,\label{asym:C2-errors}\\
 |F_t''(z)+\cos z|&\le\varepsilon_2:=(2+Z)m_1/L+m_2/L^2,
 \qquad |z|\le Z.\nonumber
\end{align}
These follow from real sine/cosine Lipschitz bounds and
$|(U/L)^2-1|\le2|U/L-1|+(U/L-1)^2$.
They are finite inequalities involving complete-source moments;
\eqref{asym:moments} does not supply a certified numerical onset time.

\begin{corollary}[Fixed positive indices]\label{asym:fixed-roots}
For each fixed $k\ge1$, the $k$th positive real zero exists and is simple
for all sufficiently large $t$, and
\begin{equation}\label{asym:root-limit}
 L_t x_k(t)\longrightarrow(k-\tfrac12)\pi.
\end{equation}
\end{corollary}
\begin{proof}
Choose disjoint small real neighborhoods of the first finitely many
positive cosine zeros.  Uniform convergence gives opposite signs at
each pair of endpoints; uniform derivative convergence gives a fixed
nonzero derivative sign inside each neighborhood.  There is exactly
one simple zero in each.  On the compact intervening intervals,
including a neighborhood of zero, uniform convergence excludes further
zeros.  These are therefore the first positive zeros in their stated
order.  No assumption about all the other complex zeros is used.
\end{proof}

For globally continued adjacent walls identified with fixed indices
$k,k+1$, this yields separation $\pi/L_t+o(L_t^{-1})$.
Their identification is an additional continuation premise.
Equation~\eqref{asym:root-limit} proves convergence of locations to zero,
not eventual monotonic root motion or a finite collision.
The distinction between two jet scales is essential:
\begin{align}
 F_t+\ii F_t'&\longrightarrow e^{-\ii z},\label{asym:rescaled-jet}\\
 P_t(z):=\frac{H(t,z/L)+\ii H_x(t,z/L)}{H(t,0)}
 &=F_t(z)+\ii L F_t'(z).\label{asym:physical-jet}
\end{align}
The second is the physical jet divided by a positive scalar.
For $M_L(z)=\cos z-\ii L\sin z$, \eqref{asym:C2-errors} gives
\begin{equation}\label{asym:jet-error}
 |P_t-M_L|\le\mathcal E:=
 \sqrt{\varepsilon_0^2+L^2\varepsilon_1^2}
 =O(\kappa^{-1/2})
\end{equation}
on a fixed real compact set.  In particular, uniformly for bounded $w$,
\begin{equation}\label{asym:interior-profile}
 P_t(k\pi+w/L)\longrightarrow(-1)^k(1-\ii w).
\end{equation}
The $C^2$ estimates also give convergence of its $w$ derivative.
The limiting phase density is $-1/(1+w^2)$: the physical phase layer
has width $O(L^{-1})$ in $z$, or $O(L^{-2})$ in the original $x$.

\subsection{Finite phase budgets for the interior and endpoint strips}
All phase changes in the next proposition are signed net changes.
They are not automatically bounds for total variation.
\begin{proposition}[A paid interior half-turn]\label{asym:phase-budget}
Fix an integer $k$ and $0<\beta<\pi/4$, and suppose
\[
 [k\pi-\pi/2-\beta,k\pi+\pi/2+\beta]\subset[-Z,Z],
 \quad L\ge1,\quad\mathcal E<1,
 \quad\varepsilon_0<\sin\beta,\quad\varepsilon_1<\cos\beta.
\]
There are adjacent simple roots $a_-,a_+$ within $\beta$ of
$k\pi-\pi/2,k\pi+\pi/2$, respectively.  For
$0<A/L<\pi/2-\beta$, let $C$ be the net phase decrease of $P_t$
between $k\pi-A/L$ and $k\pi+A/L$.  Put
\[
 r_A=\sqrt{\cos^2(A/L)+L^2\sin^2(A/L)}.
\]
Then
\begin{align}
 C_0&:=2\arctan[L\tan(A/L)],&
 |C-C_0|&\le2\arcsin(\mathcal E/r_A),\label{asym:central-bound}\\
 |\pi-C|&\le2/A+2\arcsin(\mathcal E/r_A).\label{asym:outside-bound}
\end{align}
For the two endpoint pieces ending at the cuts
$k\pi\pm(\pi/2-\beta)$, their combined net phase decrease $S$ obeys
\begin{equation}\label{asym:endpoint-bound}
 |S|\le\frac{2\tan\beta}{L}
       +2\arcsin(\mathcal E/r_\beta),\qquad
 r_\beta=\sqrt{\sin^2\beta+L^2\cos^2\beta}\ge1.
\end{equation}
The denominator $r_\beta$ may be replaced by $L\cos\beta$ only when
$\mathcal E<L\cos\beta$, so that the resulting arcsine is defined.
\end{proposition}
\begin{proof}
At the ends of each root neighborhood the cosine values have opposite
signs of magnitude $\sin\beta$.  Inside, its derivative has a fixed
sign and magnitude at least $\cos\beta$.  The hypotheses and
\eqref{asym:C2-errors} therefore give one simple actual root in each
neighborhood.  In between, $|\cos z|\ge\sin\beta$ excludes any
additional root.  The two roots are adjacent and the selected cuts
are strictly inside their lobe.  The real component $F_t$ has constant
nonzero sign there; its endpoint derivatives have opposite signs.
A continuous half-plane argument therefore gives total cap decrease
exactly $\pi$.

Since $|M_L|\ge1$, the quotient $P_t/M_L$ stays in the disk centered
at $1$ with radius $\mathcal E/|M_L|<1$.  Its principal argument defines
one continuous error function $e(z)$ on the entire interval, with
$|e(z)|\le\arcsin(\mathcal E/|M_L(z)|)$.  Subtracting its values at
the two central cuts proves \eqref{asym:central-bound}.  Moreover
\[
 \pi-C_0=2\arctan\frac1{L\tan(A/L)}\le2/A,
\]
which proves \eqref{asym:outside-bound} using the exact full decrease.
For the endpoint pieces the model decrease is
$2\arctan(\tan\beta/L)\le2\tan\beta/L$.
Each actual root has the same imaginary-axis phase as its corresponding
model root, by the derivative signs already proved.  Only the two cut
arguments need perturbation, each with model modulus $r_\beta$.
This proves \eqref{asym:endpoint-bound}.
\end{proof}

A sufficient condition to turn these net bounds into positive phase
mass bounds is
\begin{equation}\label{asym:monotone-phase}
 2\varepsilon_1+\varepsilon_0+\varepsilon_2+
       \varepsilon_0\varepsilon_2<1.
\end{equation}
Indeed direct expansion bounds $F_t'^2-F_tF_t''$ below by one minus
the left side, and
\[
 -\frac{d\arg P_t}{dz}
 =\frac{L(F_t'^2-F_tF_t'')}{F_t^2+L^2F_t'^2}>0.
\]
For fixed $Z$, the moment asymptotics imply this eventually, without
an effective onset claimed here.  Thus theta's fixed-index large-time
phase is concentrated at the interior extremum; its fixed-width scaled
endpoint strips carry vanishing phase.  These statements preserve the
ordinary half-turn and give no surplus winding.

None of the estimates is asserted uniformly for growing index $k$ or
growing compact radius $Z$.  Such a regime requires new uniform source
bounds and wall identification.  A sign at a fixed physical location
cannot silently be treated as a sign inside the moving fixed-index cap,
and spatial coalescence as $t\to\infty$ is not a boundary collision at
any finite time.


\section{Concentration, energy and invariant information}
\label{sec:information}
\subsection{Changing the derivative scale redistributes phase}
For a positive constant \(\beta\) on a cap, define
\[
 J_\beta=f+\ii\beta f',\qquad
 -d\arg J_\beta=\frac{\beta[(f')^2-ff'']}
                     {f^2+\beta^2(f')^2}\,dx.
\]
Positive scaling of the imaginary component is a homotopy through
nonzero jets, preserving winding and the root-endpoint phase increment.
The simple-root linear model has density
\(\beta/(d^2+\beta^2)\) and window mass \(\arctan(V/\beta)\).
For the quadratic, choosing \(\beta=r/2\) removes endpoint concentration
in the \(z\) coordinate. For the theta saddle cap, \(\beta=1/L_t\)
gives \(F_t+\ii F_t'\to e^{-\ii z}\), spreading phase across the lobe.
Time-dependent positive scaling preserves finite contour degree;
an \(x\)-dependent scaling adds a derivative term to the displayed density.
Thus concentration describes a representation of the phase. The
invariant target is the fully accounted surplus \eqref{eq:surplus}.

\subsection{Removing the denominator forgets the cap count}
For a root-to-root cap, integration by parts gives
\[
 \int_u^v Q\,dx=2\int_u^v(f')^2\,dx>0,\qquad
 \int_u^v\frac{Q}{f^2+(f')^2}\,dx=(N+1)\pi.
\]
Positivity of the first integral is compatible with every cap count.
For odd \(n\), the cap functions \(f_n(x)=-\sin(nx)/n\) on \([0,\pi]\)
have identical endpoint slopes \((-1,+1)\), identical \(Q\equiv1\)
and \(\int Q=\pi\), but \(N=n-1\) and \(\int Q/E=n\pi\).
These isolate lost cap information; they are not ordinary-birth
heat-flow controls.

\subsection{Other exact simplifications have different targets}
A nonzero level arc \(H/h_2=e>0\) makes the curved-side phase depend
only on its endpoints. If that arc and a cap enclose the initial
event, and the cap has \(N_{\rm level}\) simple interior roots,
the endpoint terms cancel and the physical winding is
\(-N_{\rm level}/2\). The original event already requires
\(N_{\rm level}=2\); an additional event requires \(N_{\rm level}\ge4\).
This count differs from \(N\) in the root-endpoint cap. Arc continuation,
its two intended intersections and jet nonvanishing are additional
premises. See
\href{../notes/06_HEAT_FLOW_CONTOURS_AND_GAUSSIAN_CAP_ENERGY_20261020.md}{Note 06}
and \href{../notes/08_LOCALIZED_COMPARATOR_AND_LEVEL_CONTROL_20261020.md}{Note 08}.

For \(G_s(y)=(4\pi s)^{-1/2}e^{-y^2/(4s)}\) and
\(f(y)=H(T+s,a+y)\), Gaussian convolution gives
\(\int G_s f^{(j)}=h_j\). The Hermite expansion yields
\[
 D_s:=s\int G_s(f')^2-\int G_s f^2
 =-h_0^2-sh_1^2+
 \sum_{n\ge3}\frac{n-2}{2}\frac{(2s)^n}{n!}h_n^2.
\]
At a multiple zero \(D_s\ge0\); a complete negative estimate would
exclude that initial event directly. It is a distinct target from an
extra cap winding. The full derivation and source kernel are in Note 06
and the preserved earlier draft.

The completed bounded HFCE comparison already excluded one additional
small box by a direct first-jet bound. Its whole-cap independent-range
Gaussian estimate failed to obtain a strict negative upper bound.
A separately fixed localized Laguerre average was proved positive on
its selected window. These are scoped results, not contrary evidence
against the exact root-wall identity or a proof of the physical
Gaussian deficit's sign on the failed-estimator window.
The authoritative results and receipts are in
\href{../notes/07_FIXED_GAUSSIAN_COMPARISON_20261020.md}{Note 07},
Note 08 and the HFCE progress register. Detailed numerical inputs
remain there and in the archived first manuscript.

\section{An analytical continuation program}
\label{sec:program}
The new research question is whether event-conditioned complete theta
information can force a signed surplus incompatible with an adjacent
root strip. The fixed-index large-time cosine regime itself gives the
ordinary half-turn. Its successful nonvanishing approximation cannot
also be used as a surplus certificate.

The next useful work has four connected parts.
First, make the source moment and phase-layer estimates effective in a
declared finite regime, retaining their dependence on time, scaled
position and root index. Second, establish the actual wall continuation
and cut placement needed in that regime, without assuming the desired
cap count or a uniform \(|h_2|\) floor. Third, isolate a complete-source
signed term that could exceed its endpoint deficits and all remainders.
An intermediate event-conditioned regime or a justified growing-index
regime is a possible setting; neither has been established here.
Fourth, compare any proposed strict margin with the ordinary quadratic,
atomic and cosine-profile controls before treating it as theta evidence.

Numerics may check a symbolic formula, refine an analytically derived
constant, or validate a predeclared finite inequality with complete
errors. They should not select a sign by a broad scan, replace a
continuation proof, or turn pointwise asymptotics into a uniform estimate.
A failed estimate is recorded with its information class and remaining
deficit. No execution of the new program is asserted by writing it.

The task contracts and statuses are separate files:
\begin{itemize}
\item
\href{../research-overview/03_ANALYTICAL_PHASE_SURPLUS_PROGRAM_20261020.md}
{Analytical phase-surplus program} and its
\href{../research-progress/03_ANALYTICAL_PHASE_SURPLUS_PROGRESS_20261020.md}
{matching progress register};
\item the original BW09 register for the overall later-collision route;
\item the HFCE register for the completed bounded energy comparisons.
\end{itemize}
The user-requested discussion record is
\href{../brainstorm/01_CAP_AT_INFINITY_20261020.md}{Brainstorm 01}.
It includes longer bookkeeping details and the precise finite error
interfaces used here.

\subsection{Control premises and limits of the conclusion}
The quadratic and positive atomic controls satisfy the generic heat
and root-wall identities without a later event. The two-event quartic
has the required extra winding but fails the initial-threshold premise.
Higher-multiplicity controls require their own initial index.
The atomic family lacks the unchanged theta source. Euler-factor controls do not have the
unchanged normalized theta source; their threshold lower bounds are
not supplied event locations. Weighted-square and finite-prime controls
fail full reflection, and their local collisions do not supply a finite
all-real threshold. Reflected and positive-primitive controls lack the
unchanged normalized theta square expansion. No unsupplied global
root system is assigned to any of them.

The source-tail comparison behind the large-time cosine limit is itself
shared by many sources. It is useful asymptotic information, not a
distinguishing RH sign. The repository's
\href{../../../../../../../result-summaries/12_collision_controls_and_screening.md}
{control catalogue}, the
\href{../../../../../../../result-summaries/README.md}{result index}
and \href{../../../../../../../avoid/README.md}{patterns to avoid}
govern this interpretation.
No strict theta phase surplus, later theta collision, finite-height
attainment theorem, uniform all-index estimate or global upper-threshold
improvement is established.

\section*{Acknowledgments and version provenance}
Edward Baker proposed the expanding heat contour, the cap at infinity,
and the quantitative phase-concentration questions developed here.
This revision was drafted for him with substantial assistance from
OpenAI's gpt-6-astra through Codex at recorded effort ultra on
20 October 2026. Separate agents of the same model assisted derivation,
proof checking and internal review; no specialist or cross-model review
is claimed.

The manuscript is a coherent revision of the earlier consolidation,
with numerical details routed to their notes and unchanged source records.
At the user's explicit request, the exact previous TeX source is retained
as \href{../drafts/20261020_before_phase_concentration.tex}{the earlier draft}.
This is a specifically requested preservation, not a new series of
automatic snapshot folders. The concise
\href{../DRAFT_HISTORY.md}{draft history} identifies the versions and
their reviews. The repository is
\url{https://github.com/ebbaker/shifted-zeta-positivity}.

\appendix
\section{Local collision index and complete finite event count}
\label{app:index}
This appendix retains the earlier manuscript's local proof so that
the collision-counting implication used in the main text remains
self-contained.

\begin{theorem}[Local collision index]\label{thm:index}
Let \(H\) be a nontrivial real-analytic solution of \eqref{eq:heat}
near \(p=(T,a)\). If its spatial zero at \(p\) has exact order \(m\ge2\),
then \(p\) is an isolated zero of \(J\), and
\[
                         \ind_pJ=-\lfloor m/2\rfloor .
\]
\end{theorem}
\begin{proof}
Put \(h_j=\partial_x^jH(T,a)\), and define
\[
 P_m(s,z)=\sum_{k=0}^{\lfloor m/2\rfloor}
              \frac{m!(-s)^kz^{m-2k}}{k!(m-2k)!}.
\]
This polynomial satisfies the same heat equation. If
\(w(z)=e^{-z^2/4}\), comparison of coefficients in
\(w(z-2v)/w(z)=e^{zv-v^2}\) gives
\[
 P_m(1,z)w(z)=(-2)^m w^{(m)}(z).
\]
Integration by parts shows that \(P_m(1,\cdot)\) is orthogonal, against
the positive weight \(w\), to every polynomial of degree less than \(m\).
If it had fewer than \(m\) real sign-changing zeros, multiplication by
the product of the linear factors at those zeros would give a polynomial
of one constant nonzero sign whose weighted integral is zero. This is
impossible. Hence all \(m\) roots are simple and real.
Scaling gives the same conclusion for \(s>0\).
For \(s<0\), the coefficients after extracting \(z^{m\bmod2}\) are
strictly positive as a polynomial in \(z^2\), including its constant term.
There are therefore \(m\bmod2\) simple real roots.
At \(s=0\), \(P_m=z^m\). Thus
\((P_m,(P_m)_z)\) vanishes only at \((0,0)\).

Take \(K=[-1,1]\times[-b,b]\), \(b=2m\).
For \(|s|\le1\) and \(|z|\ge b\),
\[
 \frac{|P_m(s,z)-z^m|}{|z|^m}
 \le\sum_{k\ge1}(m^2/b^2)^k=\tfrac13 .
\]
In particular the two spatial walls have nonzero real part, and all
real roots on the time faces lie inside those walls.
For any real \(C^2\) function \(f\) with nonzero endpoints and \(N\)
simple interior zeros,
\begin{equation}
 \Delta\arg(f+\ii f')
 =\arctan\frac{f'(v)}{f(v)}
  -\arctan\frac{f'(u)}{f(u)}-N\pi.               \label{eq:phasezeros}
\end{equation}
Indeed \(f'/f=(x-r)^{-1}+O(1)\) at a simple zero \(r\);
the principal arctangent jumps upward by \(\pi\), whereas the lifted
argument is continuous. Subtracting \(\pi\) at each crossing proves the
formula. On the two spatial walls the arctangent endpoint differences
are the entire phase increments. The four endpoint terms cancel,
with the upper time face traversed left to right and the lower face
right to left. Thus the polynomial jet has degree
\((m\bmod2-m)/2=-\lfloor m/2\rfloor\).

It remains to justify the passage from the polynomial to \(H\), including
isolation. Its convergent Taylor series
\(H(T+\tau,a+y)=\sum a_{kj}\tau^ky^j\) satisfies
\[
 a_{kj}=\frac{(-1)^kh_{j+2k}}{k!j!}.
\]
The leading parabolic weight \(2k+j=m\) is \(cP_m\), where
\(c=h_m/m!\ne0\); all lower weights vanish.
Choose \(r_0>0\) with absolute convergence on a larger polydisc than
\(|\tau|\le r_0^2,\ |y|\le r_0b\), also for the differentiated series.
The finite tail sums
\[
 C_0=\frac1{|c|}\sum_{2k+j\ge m+1}|a_{kj}|r_0^{2k+j-m-1}b^j,\quad
 C_1=\frac1{|c|}\sum_{\substack{2k+j\ge m+1\\j\ge1}}
                 j|a_{kj}|r_0^{2k+j-m-1}b^{j-1}
\]
give, uniformly on \(K\) for \(0<r\le r_0\),
\[
 \left|\frac{H(T+r^2s,a+rz)}{cr^m}-P_m(s,z)\right|\le rC_0,\qquad
 \left|\frac{H_x(T+r^2s,a+rz)}{cr^{m-1}}-(P_m)_z(s,z)\right|
 \le rC_1.
\]
The polynomial jet has a positive minimum norm \(\delta_m\) on
\(\partial K\). For \(r(C_0+C_1)<\delta_m/2\), the straight boundary
homotopy is nonzero. Every nonzero sufficiently small \((\tau,y)\)
belongs to such a scaled boundary, by taking
\(r=\max(\sqrt{|\tau|},|y|/b)\). This proves isolation as well.
Both the source scaling and the target normalization preserve
orientation: their determinants are \(r^3>0\) and
\(1/(c^2r^{2m-1})>0\). The degrees agree.
\end{proof}

\begin{corollary}[Complete event count]\label{cor:count}
If a bounded piecewise smooth Jordan domain has \(J\ne0\) on its boundary,
then
\begin{equation}
 W_J(\Omega)=-\sum_{p\in\Omega:J(p)=0}\lfloor m(p)/2\rfloor .
                                                        \label{eq:count}
\end{equation}
\end{corollary}
\begin{proof}
Analyticity and the heat equation exclude infinite spatial order unless
the solution is identically zero. Theorem~\ref{thm:index} gives isolation;
compactness and a zero-free boundary collar give finiteness. The form in
\eqref{eq:winding} is locally \(d\arg J\), hence closed. Remove small
disjoint event neighborhoods and apply Stokes' theorem. Their positively
oriented indices sum to the outer winding.
\end{proof}

For the earlier nonreal-root assertion, apply the scaled analytic
expansion also on compact complex \(z\) sets. At \(s=-1\) the
polynomial \(P_m(-1,z)=\ii^mP_m(1,z/\ii)\) has simple imaginary roots,
with at least one nonreal pair for \(m\ge2\).
Small disjoint disks around those roots, uniform convergence, and
Rouch\'e's theorem give nonreal roots of \(H(T-r^2,\cdot)\).
This proves the assertion used in the threshold discussion.

If a domain contains an initial event of multiplicity \(m_0\), no other
event at or before its time, and later event charge
\(K_{\mathrm{later}}=\sum_{t>T}\lfloor m(p)/2\rfloor\), then
\[
 W_J=-q_0-K_{\mathrm{later}},\qquad q_0=\lfloor m_0/2\rfloor.
\]
A complete upper bound \(W_J<-q_0\) would force a later event.
An interval excluding every integer instead contradicts a proposed
winding enclosure itself; adding events does not repair such an interval.



\begin{thebibliography}{9}
\bibitem{debruijn}
N.~G. de Bruijn, \emph{The roots of trigonometric integrals},
Duke Mathematical Journal \textbf{17} (1950), 197--226.
\url{https://doi.org/10.1215/S0012-7094-50-01720-0}.
\bibitem{newman}
C.~M. Newman, \emph{Fourier transforms with only real zeros},
Proceedings of the American Mathematical Society \textbf{61} (1976),
245--251.
\url{https://sites.math.northwestern.edu/~auffing/papers/Newman.pdf}.
\bibitem{rodgerstao}
B.~Rodgers and T.~Tao, \emph{The de Bruijn--Newman constant is non-negative},
Forum of Mathematics, Pi \textbf{8} (2020), e6.
\url{https://doi.org/10.1017/fmp.2020.6}.
\end{thebibliography}
\end{document}

